% Generated draft: requires mathematical and source review.
\documentclass[11pt]{article}
\usepackage[margin=25mm]{geometry}
\usepackage{fontspec}
\setmainfont{Latin Modern Roman}
\usepackage{amsmath,amssymb,mathtools,unicode-math}
\setmathfont{Latin Modern Math}
\usepackage{enumitem,needspace,xurl,hyperref}
\setlength{\parindent}{0pt}
\setlength{\parskip}{5pt}
\setlength{\emergencystretch}{3em}
\newcommand{\N}{\mathbb{N}}
\newcommand{\Z}{\mathbb{Z}}
\newcommand{\Q}{\mathbb{Q}}
\newcommand{\R}{\mathbb{R}}
\newcommand{\C}{\mathbb{C}}
\newcommand{\F}{\mathbb{F}}
\newcommand{\A}{\mathbb{A}}
\newcommand{\PP}{\mathbb{P}}
\newcommand{\E}{\mathbb{E}}
\newcommand{\eps}{\varepsilon}
\newcommand{\dd}{\,\mathrm{d}}
\begin{document}
\section*{UnsolvedMath ID: 800727}
\textbf{Source problem code:} catalogue entry
\subsection{Definitions and mathematical statement}
For a topological space \(X\), its weight is
\[
w(X)=\min\{\lvert\mathcal B\rvert:\mathcal B\text{ is a base for the topology of }X\}.
\]
A topological group is a group equipped with a topology for which multiplication and inversion are continuous. Fix a separation convention for topological groups and apply it consistently to all groups below. The supplied excerpt does not specify whether Hausdorffness is required, so the intended scope of this convention remains unresolved.

For this draft, a space is called zero-dimensional if it has a base consisting of sets that are both open and closed. A space \(X\) is called almost metrizable if there exist a nonempty compact subset \(K\subseteq X\) and open sets \(U_n\subseteq X\), \(n\in\N\), such that \(K\subseteq U_n\) for every \(n\), and for every open set \(U\) containing \(K\), some \(U_n\) satisfies \(U_n\subseteq U\). These are working definitions; the supplied excerpt does not give its definitions of zero-dimensionality or almost metrizability, and their agreement with the source requires review.

A quotient homomorphism \(q:H\to G\) means a continuous surjective group homomorphism such that, for every subset \(V\subseteq G\),
\[
V\text{ is open in }G\quad\Longleftrightarrow\quad q^{-1}(V)\text{ is open in }H.
\]
Thus \(G\) carries the quotient topology, in addition to being an algebraic quotient of \(H\).

\begin{enumerate}
\item Is it true that for every topological group \(G\), there exist a zero-dimensional topological group \(H\) and a quotient homomorphism \(q:H\to G\) such that
\[
w(H)=w(G)?
\]
\item Characterize the almost metrizable topological groups \(G\) for which there exist a zero-dimensional topological group \(H\) and a quotient homomorphism \(q:H\to G\) satisfying \(w(H)=w(G)\).
\end{enumerate}
In the second question, the supplied statement does not require \(H\) itself to be almost metrizable. The characterization requested there, and the universal assertion in the first question, are subject to the unresolved terminology and separation conventions stated above.
\subsection{Short English statement}
Does every topological group arise as a quotient of a zero-dimensional topological group with exactly the same weight? Which almost metrizable topological groups admit such a representation? The source's separation and terminology conventions require confirmation.
\subsection{Sources}
\begin{itemize}
\item Local UnsolvedMath registry: permanent ID \(800727\), source problem code \texttt{catalogue entry}, supplied title \emph{Open Problems in Topological Algebra --- Problem 1 (Choban)}. The record lists Banakh and Choban as proposers, with proposed year 2012, and specifically attributes the problem to Choban in its background.
\item Banakh and Choban, \emph{Some Open Problems in Topological Algebra} (2012), Problem 1. Supplied source URL: \url{https://arxiv.org/abs/1202.1619}. The record identifies the extraction method as \texttt{arxiv\_latex} and reports access on 2026-10-16; this access information is supplied provenance and has not been independently verified here.
\item The local record labels the problem \emph{open}. This is the supplied registry classification, not an independently verified assessment of the literature or current mathematical status.
\end{itemize}
\end{document}
